\section{\'Etat d'avancement du projet}
\label{sec:etat}

Afin de rester fid\`ele \`a l'\'etat r\'eel du d\'ep\^ot au moment de la remise, le tableau
ci-dessous distingue les composants termin\'es, ceux en cours de finalisation et ceux encore
\`a compl\'eter. Le suivi d\'etaill\'e est \\'egalement consign\'e dans le fichier
\texttt{PHASES.md} du d\'ep\^ot GitHub.

\begin{table}[H]
  \centering
  \caption{\'Etat des phases du projet (juin 2026).}
  \label{tab:phases}
  \begin{tabular}{@{}p{0.42\textwidth}p{0.18\textwidth}p{0.32\textwidth}@{}}
    \toprule
    \textbf{Composant} & \textbf{Statut} & \textbf{Livrable principal} \\
    \midrule
    Setup et collecte du corpus & Termin\'e & \texttt{collect\_corpus.py}, \texttt{data/raw/} \\
    Nettoyage et chunking & Termin\'e & \texttt{src/prepare\_data.py}, 1443 chunks \\
    Index vectoriel ChromaDB & Termin\'e & \texttt{src/build\_index.py}, \texttt{chroma\_db/} \\
    Pipeline RAG (CLI) & Termin\'e & \texttt{src/rag\_pipeline.py} \\
    Interface Streamlit & Impl\'ement\'ee, tests finaux en cours & \texttt{app.py} \\
    Jeu d'\'evaluation (34 questions) & Termin\'e & \texttt{evaluation/eval\_questions.json} \\
    M\'etriques et r\'esultats CSV & Termin\'e & \texttt{evaluation/evaluate.py}, \texttt{results.csv} \\
    Visualisations & Termin\'e & \texttt{visualizations/plot\_results.py}, figures PNG \\
    Rapport PDF (LaTeX) & Remis (ce document) & \texttt{report/overleaf/} \\
    Revue finale et d\'emo orale & En cours & README, script de d\'emonstration \\
    \bottomrule
  \end{tabular}
\end{table}

\subsection{Ce qui fonctionne aujourd'hui}

\begin{itemize}[noitemsep]
  \item T\'el\'echargement automatique de 36 articles Wikipedia, nettoyage et production de
        1443 chunks index\'es dans ChromaDB.
  \item Interrogation du syst\`eme en ligne de commande
        (\texttt{python -m src.rag\_pipeline}) avec affichage de la r\'eponse et des sources.
  \item \'Evaluation automatique sur 34 questions avec calcul de Precision@3, Recall@3,
        fid\'elit\'e lexicale et temps de r\'eponse ; r\'esultats stock\'es dans
        \texttt{evaluation/results.csv}.
  \item G\'en\'eration de graphiques \`a partir des r\'esultats d'\'evaluation.
\end{itemize}

\subsection{Travail restant}

\begin{itemize}[noitemsep]
  \item \textbf{Interface Streamlit} : l'application (\texttt{app.py}) propose trois onglets
        (Chat, \'Evaluation, Gestion du corpus). Les tests utilisateur finaux et la pr\'eparation
        de la d\'emonstration en direct (15 juin 2026) sont en cours.
  \item \textbf{Documentation et d\'emo} : mise \`a jour du \texttt{README.md} avec les
        instructions compl\`etes (Streamlit, \'evaluation, rapport) et pr\'eparation d'un sc\'enario
        de d\'emonstration oral.
  \item \textbf{Am\'eliorations possibles} (hors remise imm\'ediate) : re-ranking des chunks,
        m\'etrique de fid\'elit\'e plus fine (LLM-as-judge), extension du corpus.
\end{itemize}

Ce positionnement transparent permet de pr\'esenter des r\'esultats d'\'evaluation r\'eels sur le
pipeline principal, tout en indiquant clairement les parties encore en cours de stabilisation
avant la soutenance.
